\section{\texorpdfstring{\code{torch.distributed}}{torch.distributed}：把 collective 写成代码}

\subsection{初始化和 process group}

官方讲义用 \code{torch.multiprocessing.spawn} 启动多个进程，每个进程运行同一个函数，但传入不同 \code{rank}。然后用 \code{dist.init\_process\_group} 建立通信环境。若有 GPU，就用 NCCL；否则用 Gloo。

\begin{lstlisting}[caption={Lecture 7 中分布式初始化的最小结构}]
def setup(rank: int, world_size: int):
    os.environ["MASTER_ADDR"] = "localhost"
    os.environ["MASTER_PORT"] = "15623"

    if torch.cuda.is_available():
        dist.init_process_group("nccl", rank=rank, world_size=world_size)
    else:
        dist.init_process_group("gloo", rank=rank, world_size=world_size)
\end{lstlisting}

\begin{knowledgebox}{一个 rank 通常对应一个进程}
在 PyTorch DDP 风格中，常见做法是 one process per GPU。这样每个进程只管理一张卡，CUDA context、随机种子、通信 rank 都更清晰。多机训练时还要区分 local rank 和 global rank：local rank 决定本机第几张 GPU，global rank 决定全局通信身份。
\end{knowledgebox}

\subsection{all-reduce 示例}

官方代码中，每个 rank 先构造一个不同的向量：

\begin{lstlisting}[caption={all-reduce 会原地修改每个 rank 的 tensor}]
data = tensor([0., 1, 2, 3], device=cuda_if_available(rank)) + rank
dist.all_reduce(tensor=data, op=dist.ReduceOp.SUM, async_op=False)
\end{lstlisting}

若 world size 为 4，四个输入分别是 \([0,1,2,3]\)、\([1,2,3,4]\)、\([2,3,4,5]\)、\([3,4,5,6]\)，sum all-reduce 后每个 rank 都得到 \([6,10,14,18]\)。这正是 DDP 梯度同步的形状：每张卡最初有自己的梯度，通信后每张卡有相同平均梯度。

\subsection{reduce-scatter + all-gather 示例}

\begin{lstlisting}[caption={reduce-scatter 后接 all-gather 复原 all-reduce 语义}]
input = torch.arange(world_size, dtype=torch.float32,
                     device=cuda_if_available(rank)) + rank
output = torch.empty(1, device=cuda_if_available(rank))
dist.reduce_scatter_tensor(output=output, input=input,
                           op=dist.ReduceOp.SUM, async_op=False)

input = output
output = torch.empty(world_size, device=cuda_if_available(rank))
dist.all_gather_into_tensor(output_tensor=output,
                            input_tensor=input, async_op=False)
\end{lstlisting}

\begin{importantbox}{这段代码为什么是后续课程的桥}
如果做完 reduce-scatter 后就停止，每个 rank 只保留归约结果的一片；如果再做 all-gather，每个 rank 又得到完整归约结果。ZeRO/FSDP 的很多设计就在这两个点之间选择：什么时候需要完整参数，什么时候只保存分片状态。
\end{importantbox}

\subsection{barrier、同步和 deadlock}

\code{dist.barrier()} 让所有 rank 等到同一个同步点。它常用于 benchmark、调试打印、确保 warmup 完成。它也会暴露控制流不一致：如果某个 rank 进入了 barrier，另一个 rank 因条件分支跳过，程序就会卡住。

\begin{warningbox}{分布式 deadlock 常来自不一致}
常见错误包括：某个 rank 少调用一次 collective、不同 rank 的 tensor shape 不一致、某些 rank 提前异常退出、print 或 dataloader 导致顺序错位。调试时先确认所有 rank 的 collective 顺序完全一致，再看性能。
\end{warningbox}

\subsection{本章小结}

\code{torch.distributed} 的 API 看起来只是函数调用，但每个调用背后都是多进程同步协议。写分布式训练代码时，正确性约束包括数学结果、shape/dtype、调用顺序和进程生命周期。

